\section{CAPÍTULO 3: CONCLUSIONES Y RECOMENDACIONES}

\subsection{3.1. CONCLUSIONES}

\begin{enumerate}
    \item El desarrollo del proyecto ParkSys, durante las Experiencias Formativas en Situaciones Reales de Trabajo IV, permitió consolidar de manera práctica los conocimientos adquiridos en el módulo de Implementación de Pruebas de Software, evidenciando que la aplicación sistemática de pruebas unitarias, de integración y de aceptación (48 pruebas automatizadas con una tasa de éxito del 100 por ciento) constituye un factor determinante para garantizar la confiabilidad de un sistema antes de su despliegue en un entorno de producción real.
    \item La experiencia formativa permitió comprobar que el proceso de optimización del rendimiento —desde la identificación de cuellos de botella hasta la implementación de mejoras concretas, tales como la indexación de campos de alta consulta y la adopción de Lazy Loading en el frontend— exige no solo dominio técnico, sino también una metodología ordenada de diagnóstico, corrección y verificación, la cual fue aplicada de forma iterativa a lo largo del desarrollo del sistema.
    \item El trabajo colaborativo realizado durante el desarrollo, las pruebas y el despliegue de ParkSys permitió fortalecer competencias transversales relacionadas con la gestión de proyectos de software en equipo, la documentación técnica rigurosa y la comunicación efectiva con el usuario final, competencias que resultan indispensables para el desempeño profesional del futuro Ingeniero de Sistemas de Información.
\end{enumerate}

\subsection{3.2. RECOMENDACIONES}

\begin{enumerate}
    \item Se recomienda que la escuela profesional destine mayor tiempo curricular a la etapa de pruebas de software dentro del plan de Experiencias Formativas, dado que, durante el desarrollo del presente proyecto, se evidenció que la elaboración de una suite de pruebas completa (unitarias, de integración y de aceptación) demanda una carga de trabajo considerable que en ocasiones compite con los plazos establecidos para las demás etapas del proyecto.
    \item Se sugiere fortalecer, dentro de los cursos previos a las Experiencias Formativas, la enseñanza de herramientas de control de versiones y de despliegue continuo (CI/CD), ya que su dominio resultó fundamental para la entrega oportuna del sistema y podría reforzarse con mayor práctica guiada antes de enfrentar un proyecto real.
    \item Se recomienda a la coordinación académica establecer espacios de asesoría periódica más frecuentes durante la etapa de implementación de pruebas y despliegue, con la finalidad de absolver oportunamente las dudas técnicas que surgen durante estas fases y optimizar así los tiempos de ejecución del proyecto.
\end{enumerate}
